\documentclass[12pt,a4paper]{article}
\usepackage[utf8]{inputenc}
\usepackage[T1]{fontenc}
\usepackage{amsmath}
\usepackage{booktabs}
\usepackage{geometry}
\usepackage{hyperref}
\usepackage{cite}
\geometry{a4paper,top=2.5cm,bottom=2.5cm,left=2.5cm,right=2.5cm}
\hypersetup{colorlinks=true,linkcolor=blue,citecolor=blue,urlcolor=blue}

\title{\textbf{Multi-Electron Capture Enhancement via Heavy Neutral Elements:\\
Scaling Proton--Electron Capture Reactions through\\
High Atomic Number Target Media}}
\author{Eymen Aydoğan}
\date{May 19, 2026}

\begin{document}
\maketitle

\begin{abstract}
We propose that neutral (non-ionized) heavy elements with high
atomic number $Z$ provide superior target media for the PECCNAL
proton--electron capture mechanism. A neutral atom with $Z$ electrons
presents $Z$ sequential capture targets to incoming protons, with
each electron remaining available until captured. Unlike hydrogen
($Z=1$, single capture), heavy neutral atoms such as iron ($Z=26$),
lead ($Z=82$), or bismuth ($Z=83$) can sustain up to $Z$ sequential
$p + e^- \rightarrow n$ reactions before complete ionization.
The reaction compactness and high electron density of heavy elements
dramatically increases neutron production per unit volume.
\end{abstract}

\noindent\textbf{Keywords:} heavy elements, multi-electron capture,
high-Z target, neutron production, atomic number scaling

\hrule\vspace{1em}

\section{Introduction}
In the original PECCNAL framework~\cite{aydogan2026a}, hydrogen
($Z=1$) serves as the reaction medium with one electron per atom.
Here we propose that neutral heavy elements offer $Z$ times more
capture targets per atom, dramatically increasing volumetric
neutron production.

The key requirement is that the element remain \textbf{neutral}
(non-ionized), ensuring each electron is electrostatically
attracted to incoming protons.

\section{Multi-Electron Capture Sequence}

For a neutral atom with atomic number $Z$:
\begin{align}
p + e^-_{(1)} &\rightarrow n_1 + \nu_e \tag{H1}\\
p + e^-_{(2)} &\rightarrow n_2 + \nu_e \tag{H2}\\
&\vdots \notag\\
p + e^-_{(Z)} &\rightarrow n_Z + \nu_e \tag{HZ}
\end{align}

Each step reduces the atom's electron count by 1, with the
atom progressing from neutral to fully ionized ($Z^+$) after
$Z$ captures.

\section{Reaction Rate Scaling}

\begin{table}[h!]
\centering
\caption{Neutron production capacity by target element}
\vspace{0.5em}
\begin{tabular}{lcccc}
\toprule
\textbf{Element} & \textbf{Z} & \textbf{Max n/atom} & \textbf{Cost} & \textbf{Availability} \\
\midrule
Hydrogen (H) & 1 & 1 & Low & Abundant \\
Iron (Fe) & 26 & 26 & Low & Abundant \\
Lead (Pb) & 82 & 82 & Low & Abundant \\
Bismuth (Bi) & 83 & 83 & Low & Moderate \\
Uranium (U) & 92 & 92$^*$ & High & Restricted \\
\bottomrule
\multicolumn{5}{l}{$^*$ Uranium also undergoes fission upon neutron capture: $n + {}^{235}\text{U} \rightarrow 200\,\text{MeV}$}
\end{tabular}
\end{table}

\section{Volumetric Neutron Production}

For iron ($Z = 26$) vs hydrogen ($Z = 1$) at the same atomic density:
\begin{equation}
\frac{\dot{n}_\text{Fe}}{\dot{n}_\text{H}}
= Z_\text{Fe} \times \frac{n_\text{Fe}}{n_\text{H}}
= 26 \times \frac{\rho_\text{Fe}/M_\text{Fe}}{\rho_\text{H}/M_\text{H}}
= 26 \times \frac{7874/55.85}{0.09/1.008}
\approx 26 \times 1590 \approx 41,000\times
\end{equation}

Iron provides approximately 41,000 times more neutrons per unit
volume than hydrogen gas at STP.

\section{Coulomb Barrier Analysis}

For neutral atoms, the incoming proton experiences attraction
from the outer electron cloud before reaching the nucleus:
\begin{equation}
U(r) = -\frac{e^2}{4\pi\epsilon_0 r}
+ \frac{Ze^2}{4\pi\epsilon_0 r}
\quad \text{(simplified, outer electron dominates)}
\end{equation}

At large distances ($r > r_\text{atom}$), the atom appears neutral
and the proton experiences no net force. As it enters the electron
cloud, it is attracted to the nearest electron.
\textbf{No Coulomb barrier exists for proton--electron interaction.}

\section{Practical Considerations}

Lead (Pb, $Z=82$) is proposed as optimal:
\begin{itemize}
\item Low cost, widely available
\item Already used as radiation shielding
\item High electron density: $n_e = 3.3 \times 10^{30}$ m$^{-3}$
\item Stable isotopes, no fission risk
\item Liquid at moderate temperatures (melting point 327°C)
\end{itemize}

\section{Conclusion}
Neutral heavy elements with high atomic number $Z$ multiply
neutron production capacity by up to $Z \times 10^3$ compared
to hydrogen gas. Lead ($Z=82$) is proposed as the optimal
practical target, providing high electron density without
fission risk. This scaling principle enables compact,
high-output PECCNAL reactor designs.

\begin{thebibliography}{9}
\bibitem{aydogan2026a} E. Aydoğan, ``Coulomb-Barrier-Free Neutron Production,'' arXiv preprint, 2026.
\bibitem{pdg2022} R. L. Workman et al. (PDG), \textit{Prog. Theor. Exp. Phys.}, 083C01, 2022.
\end{thebibliography}
\end{document}
